\section{Methods}
\label{sec:methods}

\begin{figure*}[t]
\centering
\includegraphics[width=\textwidth]{figures/fig1_workflow.pdf}
\caption{\rev{Study workflow. For each held-out participant $p$, z-scoring, model fitting and channel ranking use the other participants $R_{-p}$; attributions are computed on participant-disjoint inner folds of $R_{-p}$; $k$-electrode subsets are retrained and tested on $p$. The same workflow was applied separately to each dataset.}}
\label{fig:workflow}
\end{figure*}

\subsection{Datasets}

\paragraph{GRABMyo (primary)}
GRABMyo v1.1.0~\cite{pradhan2022multiday} contains recordings from \GRABMyoParticipants\ able-bodied participants (\GRABMyoMale\ male, \GRABMyoFemale\ female; age \GRABMyoAgeMean\ $\pm$ \GRABMyoAgeSD\ years) on \GRABMyoSessions\ different days. Signals were recorded with an EMGUSB2+ amplifier (OT Bioelettronica) at 2048~Hz. In each session participants performed 16 hand and wrist gestures plus rest, with 7 trials of 5~s each. We analysed the 16 forearm electrodes. These form two circumferential rings of 8 electrodes (F1--F8 and F9--F16), 2~cm apart along the forearm, with electrode $j$ of one ring aligned with electrode $j+8$ of the other. Mean circumference at the ring position was \GRABMyoCircMean\ $\pm$ \GRABMyoCircSD~cm, giving about \GRABMyoArcSpacing~cm between neighbouring electrodes in a ring. The recordings are monopolar, so the selection unit in the primary analysis is a single electrode. All \GRABMyoFilesVerified\ distributed files matched the dataset's published SHA-256 checksums. All \GRABMyoRecords\ records were processed, with \GRABMyoRecordsFailed\ failures. The primary task is the 16 active gestures (rest excluded), for which chance Macro-F1 is about 0.063.

\paragraph{NinaPro DB5 (exploratory methodological replication)}
DB5~\cite{pizzolato2017comparison} provides \DBFiveSubjects\ able-bodied subjects (\DBFiveMale\ male, \DBFiveFemale\ female; age \DBFiveAgeMean\ $\pm$ \DBFiveAgeSD\ years, from the files' metadata) recorded with two Myo armbands (16 channels, 200~Hz). We used exercise~1 (12 finger movements, 6 repetitions each), segmented with the refined \texttt{restimulus}/\texttt{rerepetition} labels. The DB5 values are 8-bit integer codes, not physical units. The spectra show no powerline peak, and their power extends to the Nyquist frequency. DB5 was analysed separately and never pooled with GRABMyo. Its channel numbers bear no anatomical correspondence to GRABMyo's, so only the behaviour of the procedure is compared across datasets.

\rev{\paragraph{putEMG (exploratory methodological replication, added post hoc)}
putEMG~\cite{kaczmarek2019putemg} provides 44 able-bodied participants recorded on two days at least one week apart (7 to 82 days in these files; median 14). A 24-electrode monopolar array forms three circumferential bands of 8 electrodes (1--8, 9--16, 17--24) on the forearm, sampled at 5120~Hz with a 12-bit converter. Each session contains three trajectories with 20 repetitions of each of 7 active gestures, labelled from video-verified ground truth. This analysis was added at the end of the study, after all GRABMyo and DB5 results existed; its design was recorded before any putEMG model was fitted, and it is exploratory. All 264 gesture files were readable and complete, and every participant had both sessions. As for DB5, the putEMG channel numbers have no correspondence with GRABMyo's.}

\subsection{Preprocessing and features}
GRABMyo signals were band-pass filtered (20--450~Hz, 4th-order zero-phase Butterworth) and notch filtered at 50~Hz ($Q = 30$). A spectral audit chose 50~Hz over 60~Hz: mean band-power ratio was 10.35 at 50~Hz and 0.62 at 60~Hz. Windows were 250~ms with 50\% overlap and never crossed trial boundaries. From each channel we extracted ten features: mean absolute value (MAV), root mean square, variance, waveform length, zero crossings (ZC), slope sign changes (SSC), integrated EMG (IEMG), Willison amplitude (WAMP), mean frequency and median frequency. This gives a 160-dimensional vector per window, \GRABMyoWindows\ windows in total and \GRABMyoActiveWindows\ from active gestures. The ZC, SSC and WAMP amplitude threshold was set to $10^{-2}$~mV, about 1.15 times the median absolute first difference of the \rev{unfiltered signal in 54 probe records from 6 participants (rounded to one significant figure). This label-free scalar was set once rather than per fold. To check whether this let a held-out participant influence the features, we recomputed the median in every outer fold without that participant (post hoc). The same rule gave $10^{-2}$~mV in \LOOMonoIdentical\ of \LOOMonoFolds\ folds (leave-one-out medians \LOOMonoMin--\LOOMonoMax${}\times 10^{-3}$~mV; only \LOOMonoInProbe\ folds held out a probe participant). Fold-wise estimation would therefore have produced identical features.} An earlier threshold 87 times smaller had saturated WAMP and was discarded. For DB5 we used a 10--95~Hz band-pass, no notch, 250~ms windows (50 samples) and a threshold of 7 ADC units, derived in the same way \rev{from all ten subjects; the leave-one-out recomputation gave the same value in \LOODBFiveIdentical\ of \LOODBFiveFolds\ folds}. \rev{For putEMG, raw values were converted to mV and processed as for GRABMyo (20--450~Hz band-pass, 50~Hz notch, 250~ms windows of 1280 samples, 50\% overlap), keeping only windows that lay entirely inside one active-gesture segment (idle and rejected samples excluded). Because raw putEMG differences are quantised at one converter step, which equals their median, the threshold rule was applied to the filtered signal; it gave \PutThresh${}\times 10^{-3}$~mV and the same value in \PutThreshIdentical\ of 44 leave-one-out folds. This gives 240 features per window.}

\subsection{Classifiers and leave-one-subject-out evaluation}
All evaluation is leave-one-subject-out (LOSO): 43 folds for GRABMyo and 10 for DB5. In fold $p$, z-scoring statistics, model fitting and channel ranking use only $R_{-p}$. \rev{Apart from the label-free threshold checked above and the training-set sizes described below, the test participant contributes only to the score.} XGBoost used 300 trees, depth 6, learning rate 0.1, row and column subsampling of 0.8, and the \texttt{hist} method. It was trained on all \FoldTrainingWindows\ outer-training windows. For compute reasons, the RBF support vector machine ($C = 1$, $\gamma$ = \texttt{scale}) used a stratified 40,000-window subsample and the random forest (200 trees, minimum leaf size 2) a 200,000-window subsample. On one calibration fold, halving the SVM subsample from 40,000 to 20,000 windows lowered Macro-F1 from \SVMCalibFortyK\ to \SVMCalibTwentyK\rev{, suggesting that the SVM had not saturated; its reported result is therefore likely a lower bound. Model hyperparameters were fixed in advance and not tuned. The training-set sizes (the SVM and random-forest subsamples and the 200,000-window cap used for the channel analyses) were chosen from held-out Macro-F1 on two calibration folds (participants 1 and 7) before the confirmatory runs, a departure from strict fold-wise isolation for those two participants (Section~\ref{sec:deviations}). XGBoost used the \texttt{multi:softprob} objective and library defaults for all unlisted parameters.} The primary metric is participant-level Macro-F1 over the 16 active gestures; accuracy and balanced accuracy were also recorded. XGBoost was prespecified as the primary model for the channel analyses on grounds of compute.

\subsection{Cross-fitted channel attribution and the SR-CFA score}
Within each outer fold, $R_{-p}$ was split into $K = 5$ participant-disjoint inner folds. For inner fold $j$, an XGBoost model was trained on the other inner folds (capped at a \rev{uniformly sampled} 200,000-window subsample and z-scored with inner-training statistics only). Exact TreeSHAP attributions~\cite{lundberg2020local} were then computed for up to 20,000 windows of the inner-held-out participants\rev{, sampled uniformly without replacement from their pooled windows. The attribution is therefore window-weighted; because every participant contributes a similar number of windows (median \TestWindowsPerParticipant), this is close to participant-weighted}. Per-window channel importance averaged absolute attributions over the 16 classes and the 10 features of the channel,
\begin{equation}
\Phi_{i,c} = \frac{1}{G F}\sum_{g=1}^{G}\sum_{f=1}^{F} \lvert \phi_{i,g,c,f} \rvert ,
\end{equation}
and was averaged over windows and normalised to sum to one across channels, giving $I_{j,c}$. Across the five inner folds, $A_c = \operatorname{median}_j I_{j,c}$ measures attribution magnitude. $U_c = \operatorname{median}_j \lvert I_{j,c} - A_c \rvert$ measures its instability across disjoint groups of training participants. Channels were ranked \rev{(ties broken by lower channel index)} by
\begin{equation}
Q_c(\lambda) = A_c - \lambda\, U_c , \qquad \lambda \in \{0, 0.25, 0.5, 1\}.
\label{eq:srcfa}
\end{equation}
$\lambda = 0$ is cross-fitted SHAP magnitude; $\lambda > 0$ is the stability-regularised score (SR-CFA). A leakage assertion verified in every fold that no inner-held-out participant appeared in its inner training set. \rev{Unit tests verified that the outer test participant is never attributed, that inner folds are participant-disjoint, and that a synthetic dataset with a participant-confounded signal detects scaler leakage.}

The frozen plan required $\lambda$ to be chosen on outer-training data only but did not specify a criterion, and no selection procedure was implemented. We therefore select no $\lambda$ and report all four (a post-hoc specification, Section~\ref{sec:deviations}).

\subsection{Sparse sensor curve and non-inferiority (H1)}
For every outer fold, $\lambda$ and $k \in \{1,\ldots,16\}$, the top-$k$ channels of the fold's own ranking were used to retrain XGBoost from scratch. Each model used the fold's fixed 200,000-window stratified outer-training subsample and was tested on the held-out participant. The $k = 16$ model of the same curve is the reference, so every paired difference $\Delta_{p,k} = \mathrm{F1}_{p,k} - \mathrm{F1}_{p,16}$ compares models with identical training budgets. Non-inferiority at margin $\delta$ requires the lower bound of the 95\% participant-cluster bootstrap interval of $\bar\Delta_k$ to exceed $-\delta$. We define $k^* = \min\{k < 16 : \mathrm{LCB}_{95}(\bar\Delta_k) > -\delta\}$, recorded as ``none'' if no $k$ qualifies. \rev{The primary margin $\delta = 0.02$ is an absolute difference in class-averaged F1 on a 0--1 scale. It was fixed before the curve was computed as the largest loss we would accept in exchange for fewer electrodes. Relative to the full-array curve reference it is about \MarginPctOfRef\% of performance, and about \MarginPctAboveChance\% of the distance above chance (1/16 under uniform guessing). It is an engineering tolerance: it has not been validated against clinical outcomes or user performance. Because $k^*$ depends strongly on $\delta$, the prespecified sensitivity margins $\delta = 0.01$ and $0.03$ are reported so that readers can apply a different tolerance. The plan also required Holm correction within H1. The primary analysis above did not apply it; we therefore also report, post hoc, one-sided cluster-bootstrap $p$-values for $H_0\colon \bar\Delta_k \le -\delta$ with Holm correction over the 15 tested $k$ at $\alpha = 0.025$ (Section~\ref{sec:deviations}).}

\subsection{Comparison with other selection methods (H2)}
At $k \in \{4, 6, 8, 12\}$ we compared six rankings, each computed on outer-training data only:
\begin{itemize}
\item SR-CFA ($\lambda = 0.5$);
\item cross-fitted SHAP ($\lambda = 0$);
\item permutation importance\rev{: from a 40,000-window stratified subsample of the outer-training data, a model was fitted on 20,000 windows and each of the 160 features was permuted in 8,000 of the remaining windows (3 repeats, Macro-F1 scoring); a channel's score is the mean over its 10 features};
\item mutual information \rev{between each feature and the gesture label, estimated on the same 20,000 windows; a channel's score is the mean over its 10 features};
\item a geometric baseline, which selects electrodes by farthest-point sampling on the $2 \times 8$ montage \rev{(circumferential arc distance on an assumed 24~cm circumference, 2~cm between rings; deterministic start at F1, ties to the lower index). It uses no data};
\item random subsets: 100 draws per participant at each $k$, and 500 at the prespecified $k^*$ \rev{($k$ electrodes without replacement within a draw; draws independent, so a subset can recur; seed fixed per participant)}.
\end{itemize}
The value $\lambda = 0.5$ for this comparison had no recorded rationale. Because the P4 comparison includes $\lambda = 0$ as well, both rankings are reported. Drawing hundreds of random subsets per participant was only tractable with a lighter comparison model (XGBoost with 50 trees, trained on a 150,000-window stratified subsample), applied identically to every method. Its absolute scores are therefore lower than those of the sparse curve and are compared only within the comparison. H2 contrasts each attribution ranking with the participant's mean over its random subsets (Holm over 8 tests). We also report where the SR-CFA subset falls within each participant's random distribution. \rev{Here and below, ``SR-CFA'' refers to $\lambda = 0.5$ and ``cross-fitted SHAP'' to $\lambda = 0$.}

\subsection{Stability regularisation (H3) and faithfulness (H4)}
H3 compared each $\lambda > 0$ with $\lambda = 0$ on the sparse curve at $k \in \{4, 6, 8, 12\}$ (Holm over 12 tests). For H4 we used deletion with retraining~\cite{hooker2019benchmark}. For $n \in \{2, 4, \ldots, 14\}$ \rev{channels retained}, the model was retrained after removing either the $16 - n$ highest-ranked channels (\rev{Remove-Top}: keeping the $n$ lowest) or the $16 - n$ lowest-ranked (\rev{Remove-Bottom}: keeping the top $n$). \rev{Remove-Random} is the mean over the random subsets above, which exist for $n \in \{4, 6, 8, 12\}$. All three arms used the same comparison model and training subsample. As a check on protocol equivalence, the \rev{Remove-Bottom} arm selects the same electrodes as the H2 cross-fitted SHAP subsets in all \FaithProtocolCells\ matching participant-by-$k$ cells. Its Macro-F1 agrees with them to within a mean absolute difference of \FaithProtocolMeanDiff\ (maximum \FaithProtocolMaxDiff). The residual differences are numerical: z-scoring was applied per subset rather than once over all features. A faithful ranking gives \rev{Remove-Top $<$ Remove-Random $<$ Remove-Bottom}. Paired contrasts were Holm-corrected within H4.

\subsection{Statistics}
The participant is the statistical unit ($n = 43$; $n = 10$ for DB5\rev{; $n = 44$ for putEMG. We use ``participant'' for GRABMyo and putEMG and ``subject'' for DB5, following each dataset's documentation}). Windows are never treated as independent. Confidence intervals are 95\% percentile intervals from 10,000 cluster-bootstrap resamples of participants. Paired comparisons use the Wilcoxon signed-rank test with the matched-pairs rank-biserial correlation $r_{rb}$ as effect size. Holm correction is applied within each confirmatory family. Exploratory analyses are uncorrected and labelled as such. \rev{A non-significant difference is reported as ``no statistically detectable difference'', never as equivalence; no equivalence test was performed.} No participant was excluded.

\subsection{Exploratory analyses}
\rev{Exploratory analyses are described in the order in which their results are reported.}

\rev{\subsubsection{Ranking stability}} We measured the agreement of rankings computed from single participants' attributions, from the five inner folds, and from the 43 outer folds. \rev{A participant's profile is the mean of their normalised out-of-fold attribution over the outer folds in which they were an inner-held-out participant. Agreement between two rankings is Spearman's $\rho$ and Kendall's $\tau$ over the 16 channel ranks, and Jaccard@$k$, the size of the intersection of the two top-$k$ sets divided by the size of their union; single-participant agreement uses all 903 participant pairs.} Outer folds share about 98\% of their training participants, so their agreement is only an upper bound.

\rev{\subsubsection{Cross-session}} XGBoost (full configuration) was trained on session~1 of $R_{-p}$ (CS1) and tested on each session of $p$. It was also trained on sessions 1 and 2 (CS2) and tested on session 3 of $p$. CS2 has roughly twice the training data of CS1, so the two conditions differ in both session diversity and quantity.

\rev{\subsubsection{Bipolar montage}} To express the selection unit as a \rev{bipolar pair}, we derived $B_j = F_j - F_{j+8}$ ($j = 1, \ldots, 8$) and re-extracted features. \rev{These are derived pairs of the monopolar recording, not physical bipolar electrodes.} The count-feature threshold was re-derived for this signal ($5 \times 10^{-3}$~mV), because the bipolar signal is about half the monopolar amplitude. \rev{How this value was derived is not fully recorded: on the documented probe records the bipolar median is \LOOBipPooled${}\times 10^{-3}$~mV, for which the stated rule gives $\LOOBipRule \times 10^{-3}$~mV. The leave-one-out medians vary little, so the value is not driven by any single held-out participant, but its provenance is a limitation of this exploratory analysis.} We then repeated the LOSO baseline, the cross-fitted ranking and the sparse curve\rev{. The reference is the $k = 8$ model of the same curve (200,000-window training subsample), not the full bipolar baseline trained on all windows.}

\rev{\subsubsection{NinaPro DB5}} The LOSO baseline, cross-fitted ranking ($K = 5$) and sparse curve were repeated on DB5, using all outer-training windows and the $k = 16$ model as reference. \rev{The plan classed DB5 as exploratory, so its primary decision is uncorrected; the Holm-adjusted result is also reported.}

\rev{\subsubsection{putEMG} The same procedure was applied to putEMG: LOSO over 44 participants; XGBoost with the headline configuration trained on all outer-training windows (at most \PutTrainWindowsMax, below the 200,000-window cap); cross-fitted ranking ($K = 5$, all four $\lambda$); and the $\lambda = 0$ sparse curve for $k = 1, \ldots, 24$ with the $k = 24$ model as reference. Non-inferiority used the same bootstrap, margins and $k^*$ rule, with Holm correction over the 23 tested $k$ also reported. Ten random subsets per participant at $k \in \{6, 12, 18\}$ were compared with the ranked subsets (Wilcoxon, Holm over 3). No hyperparameter or training-set size was calibrated on putEMG. Because the montage has 24 electrodes in three bands, $k$ values are not comparable with GRABMyo's.}

\rev{\subsubsection{Fixed layout}} In a post-hoc analysis, a single 12-electrode layout was derived from a cross-fitted ranking ($\lambda = 0.5$) on the 22 odd-numbered participants. Models with 16 and with 12 electrodes were trained on those participants and tested on the 21 even-numbered participants. This is one split, not a LOSO estimate.

\rev{\subsubsection{Feature redundancy}} IEMG equals 512 times MAV for these windows. We repeated the comparison-model LOSO with IEMG removed and compared both Macro-F1 and channel attribution.

\subsection{Deviations from the analysis plan}
\label{sec:deviations}
The following changes were made after results existed. None alters a prespecified decision rule.
\begin{enumerate}
\item \rev{Holm correction within H1, required by the plan, was not implemented in the primary analysis. It was added post hoc and gives the same $k^*$ at all three margins (Supplementary Table~\ref{tab:h1_holm}). A fixed-sequence analysis (testing $k = 15, 14, \ldots$ and stopping at the first failure) is reported only as a post-hoc sensitivity analysis.}
\item \rev{Training-set sizes were calibrated on held-out folds 1 and 7 before the confirmatory runs.}
\item \rev{The leave-one-out check of the count-feature thresholds was added post hoc.}
\item $\lambda$ was not selected; see above.
\item The H2--H4 families were specified in detail (tests and $k$ sets) after the sparse curve and selection-baseline results had been computed.
\item The comparison model for H2 and H4 was reduced from 150 to 50 trees for compute reasons before any baseline result existed.
\item An earlier faithfulness run mixed arms trained with different budgets and was discarded, and all arms were recomputed under one protocol.
\item An earlier bipolar run used a count threshold inherited from the monopolar signal, and was discarded and recomputed.
\item The fixed-layout analysis was not planned.
\item The window-length sensitivity and the 17-class task, both listed as exploratory in the plan, were not run.
\item \rev{The putEMG replication was not in the plan and was added after all other results existed.}
\end{enumerate}
